%----------------------------------------------------------------------------------------
%	CHAPTER 6: CONCLUSIONS
%----------------------------------------------------------------------------------------

\chapter{Chapter 6. Conclusion and Future Work}
\label{ch:conclusions}

The main goal of this thesis was to compare compression methods for
dense retrieval on equal terms. Three encoder backbones with different
retrieval priors were fine-tuned under identical conditions, evaluated
on the same benchmark, and subjected to the same set of post-hoc,
training-time, and stacked compression techniques.

%----------------------------------------------------------------------------------------
\section{Conclusion}
\label{sec:conclusion}

\paragraph{Post-hoc versus training-time compression.}
Training-time compression is not always better than post-hoc
compression. The winner depends on the backbone. BGE was pre-trained
at a large scale with a contrastive retrieval objective, so it begins
from a retrieval-friendly representation. For BGE, post-hoc PQ is
ahead in the Pareto frontier at all compression ratios that were
tested, and it beats all training-time compression. But for RoBERTa
and MPNet, which do not have retrieval pre-training, training-time
binarization is on the frontier, especially at 32$\times$ compression.

\paragraph{STE versus annealed tanh for binary embedding training.}
For training-time binarization, the selection of a gradient surrogate
depends on the backbone. In the standalone case, STE does better for
both BGE and MPNet, and annealed tanh does better for RoBERTa. In the
stacked compression case, annealed tanh combines with MRL much more
cleanly than STE. MRL$+$BAT-atanh beats MRL$+$BAT-STE on all backbones
and matches or surpasses MRL$+$post-hoc binary.

\paragraph{Stacking complementary methods.}
When dimension reduction and precision reduction are stacked, the
frontier is significantly pushed beyond what any approach can achieve
on its own. For BGE, MRL$+$PQ outperforms standalone PQ in the range
of 48$\times$ to 96$\times$. But not all combinations are beneficial;
for instance, on BGE and MPNet, joint MRL$+$BAT training performs
worse than the sequential MRL$\rightarrow$post-hoc binary.

\paragraph{Cross-backbone generalizability.}
These results are not consistent across backbones. Below 8$\times$,
the backbone does not matter: INT8 and 4-bit TurboQuant are almost
lossless on all three, and the order is the same on all backbones
when techniques are ordered by quality.

At 32$\times$, the backbone starts to matter. While the order of
post-hoc methods remains the same (PQ $>$ TurboQuant-1bit $>$ post-hoc
binary), for training-time methods the relative order changes. For
BGE, BAT lowers the scores below the float32 baseline but pushes
RoBERTa above its baseline. Therefore, a study conducted only on a
retrieval-specialized backbone cannot be applied to general-purpose
backbones at high compression
(Section~\ref{sec:cross_backbone}).

This thesis finds that there is no single compression technique that
dominates across the compression range and across different backbones.
The right compression method depends on two things: whether the
encoder is pre-trained for retrieval, and the target memory budget.
These studies are valuable when the compression technique needs to be
selected during dense retrieval deployment.

%----------------------------------------------------------------------------------------
\section{Future Work}
\label{sec:future_work}

There are several limitations we have faced, and they point to natural
next steps.

\paragraph{Scale.}
NanoBEIR allowed for rapid performance testing with limited resources,
but each subset has a limited corpus size between 2{,}000 and 6{,}000.
We can utilize the full BEIR or MTEB benchmark to see if the result
stays the same or not. The training mixture dataset was also limited;
we can include more data points from diverse sources. Each model is
trained for only one epoch; a longer training schedule could improve
training-time compression techniques.

\paragraph{MRL with recurrent binarization.}
We could pair MRL with the recurrent residual binarization that was
deployed at Tencent~\cite{gan2023binary}. This approach encodes
embeddings in binary bits, where each additional bit corrects previous
residual errors, allowing the user to specify the desired bit depth at
inference. This allows post-hoc control over both the truncation
dimension and the bit depth.

\paragraph{Why BAT helps RoBERTa more than BGE.}
BAT raises RoBERTa above its float32 baseline but drops BGE below its
own. This is not just slow convergence of the float32 model: the gap is
there from the first checkpoint and does not close with more training
(Figure~\ref{fig:training_trajectory}). Why BAT beats the uncompressed
baseline on RoBERTa stays open. One hypothesis is that BGE's float32
embeddings already carry useful retrieval signals, so adding a sign
destroys vital information.
In RoBERTa's case, these embeddings are mostly noise, so a similar
sign transformation causes less harm. To test this hypothesis, we can
add controlled noise to BGE's weights before fine-tuning them. If our
hypothesis is right, then as the noise magnitude increases, we should
see BAT's relative quality improve. Another analysis for this could be
comparing the mean-pooled embedding before binarization of BGE and
RoBERTa.